\documentclass[11pt]{article}
\usepackage[a4paper,margin=18mm]{geometry}
\usepackage{fontspec}
\setmainfont{Latin Modern Roman}
\usepackage{amsmath,amssymb,amsthm,amscd,graphicx,color}
\usepackage{xurl}
\usepackage[unicode,hidelinks]{hyperref}
\allowdisplaybreaks
\setlength\emergencystretch{3em}
\numberwithin{equation}{section}

\newtheorem{Theorem}{Theorem}[section]
\newtheorem{Corollary}[Theorem]{Corollary}
\newtheorem{Lemma}[Theorem]{Lemma}
\newtheorem{Proposition}[Theorem]{Proposition}


\begin{document}
\begin{center}\Large Balanced terminating multivariable elliptic Gustafson-Rakha summation for arbitrary dimension and both parities\end{center}
\noindent\textbf{Stable ID:} 1883\par
\noindent\textbf{Authors:} Hjalmar Rosengren\par
\noindent\textbf{Original work:} Gustafson-Rakha-Type Elliptic Hypergeometric Series\par
\noindent\textbf{Version:} Official SIGMA 13 (2017) article 037, DOI 10.3842/SIGMA.2017.037; journal PDF and native TeX acquired 2026-10-01, exact bytes pinned by SHA256\par
\noindent\textbf{Selected task scope:} Theorem3.1 only: integer n>=\allowbreak{}1,N>=\allowbreak{}0, nonzero q,z\_i,b\_j, elliptic nome p of absolute value less than one, q\textasciicircum{}(N-1)b1b2b3b4 product(z\_i\textasciicircum{}2)=\allowbreak{}1, and parameters where all displayed denominator factors are nonzero. Retain both n parity cases of the exact author formula; meromorphic continuation is not a separate claim.\par
\noindent\textbf{Difficulty:} H2: The author proves a previously conjectured finite multivariable elliptic summation by two coupled inductions. Its core combines two-point theta interpolation with a reindexing of the composition sum and a second theta partial-fraction identity, while retaining distinct parity-dependent products. This is a substantive special-function identity rather than substituting parameters into an existing summation.\par
\noindent\textbf{Editorial boundary note (not author text):} The complete original main Theorem 3.1 is retained, including both parity-dependent products and the complete author two-part induction. The N=\allowbreak{}1 theta-interpolation seed, N-to-N+\allowbreak{}1 recurrence, composition reindexing, all q-exponent cancellations and the final inner theta sum are supplied. Exact theta/Pochhammer/Delta notation and denominator-nonvanishing parameter scope are recorded. Classical Weierstrass interpolation and theta partial fractions remain cited standard background. Later corollaries, Bailey transformations and separate meromorphic-continuation claims are not counted.\par
\noindent\textbf{Source:} \url{https://sigma-journal.com/2017/037/}\quad DOI: \url{https://doi.org/10.3842/SIGMA.2017.037}\par
\noindent\textbf{License and attribution:} Copyright retained by the named authors. Published by SIGMA under CC BY 4.0: \url{https://creativecommons.org/licenses/by/4.0/}. Source: \url{https://sigma-journal.com/about.html}. Adaptation consists of standalone excerpt formatting, cover and numbering restoration; original author language and mathematical text are retained.\par
\medskip\hrule\medskip\noindent\textbf{Author text begins below.}\par
\setcounter{section}{1}
\setcounter{subsection}{0}
\setcounter{subsubsection}{0}
\setcounter{Theorem}{0}
\setcounter{equation}{1}
% Original source lines 68--91
\section{Preliminaries}

When $z=(z_1,\dots,z_n)$ is a vector we will write $|z|=z_1+\dots+z_n$ and $Z=z_1\dotsm z_n$.

Throughout, $p$ and $q$ will be f\/ixed parameters with $|p|<1$. We employ the standard notation
\begin{gather*}\theta(z)=\prod\limits_{j=0}^\infty\big(1-p^jz\big)\big(1-p^{j+1}/z\big), \\
(z)_k=\begin{cases}\theta(z)\theta(qz)\dotsm\theta\big(q^{k-1}z\big), & k\in\mathbb Z_{\geq 0},\\
1/\theta\big(q^kz\big)\theta\big(q^{k+1}z\big)\dotsm\theta\big(q^{-1}z\big), & k\in\mathbb Z_{<0}\end{cases} \end{gather*}
as well as
\begin{gather*}\theta(z_1,\dots,z_m)=\theta(z_1)\dotsm\theta(z_m), \qquad (z_1,\dots,z_m)_k=(z_1)_k\dotsm(z_m)_k.\end{gather*}
Most of our computations are based on the elementary identities
\begin{gather*}\theta(1/z)=\theta(pz)=-\theta(z)/z, \\
(a)_{n+k}=(a)_n(aq^n)_k,\qquad (a)_{n-k}=(-1)^kq^{\binom k2}(q^{1-n}/a)^k\frac{(a)_n}{(q^{1-n}/a)_k}, \end{gather*}
which will be used without comment. All our sums contain the $A$-type factor
\begin{gather*}\frac{\Delta(zq^x)}{\Delta(z)}=\prod\limits_{1\leq i<j\leq n}\frac{q^{x_i}\theta(q^{x_j-x_i}z_j/z_i)}{\theta(z_j/z_i)}, \end{gather*}
where $z=(z_1,\dots,z_n)$ and $x=(x_1,\dots,x_n)$ is the summation index. We mention the useful identity \cite[equation~(3.8)]{r}
\begin{gather}\label{etf}
\frac{\Delta(z q^x)}{\Delta(z)}=(-1)^{|x|}q^{-\binom{|x|}2-|x|}\prod\limits_{i,j=1}^n\frac{(qz_i/z_j)_{x_i}}{(q^{-x_j}z_i/z_j)_{x_i}}
\end{gather}
and \cite[Example~20.53.3]{ww}
\begin{gather}\label{tpf}
\sum_{k=1}^n\frac{\prod\limits_{j=1}^{n+1}\theta(z_k/b_j)}{\theta(z_k/t)\prod\limits_{j=1,\, j\neq k}^n\theta(z_k/z_j)}=\frac{\prod\limits_{j=1}^{n+1}\theta(b_j/t)}{\prod\limits_{j=1}^n\theta(z_j/t)},
\end{gather}
valid for $tz_1\dotsm z_n=b_1\dotsm b_{n+1}$.

\setcounter{section}{2}
\setcounter{subsection}{0}
\setcounter{subsubsection}{0}
\setcounter{Theorem}{0}
\setcounter{equation}{5}
% Original source lines 109--190
\section{The elliptic Gustafson--Rakha summation}\label{ms}

Our main result is the following identity. It is easy to see that the case $p=0$ is equivalent to \cite[Theorem~1.2]{gr} and the general case to the conjecture of \cite[p.~953]{sp}. Recall the notation $Z=z_1\dotsm z_n$.

\begin{Theorem}\label{grt}
For parameters subject to $q^{N-1}b_1\dotsm b_4z_1^2\dotsm z_n^2=1$,
\begin{gather}\sum_{\substack{x_1,\dots,x_n\geq 0,\\x_1+\dots+x_n=N}}
\frac{\Delta(zq^x)}{\Delta(z)}\prod\limits_{1\leq i<j\leq n}q^{x_ix_j}(z_iz_j)_{x_i+x_j}\prod\limits_{i=1}^n
\frac{\prod\limits_{j=1}^4(z_ib_j)_{x_i}}{z_i^{x_i}\prod\limits_{j=1}^n(qz_i/z_j)_{x_i}}\nonumber\\
\qquad{} =\begin{cases}
\displaystyle\frac{(Zb_1,Zb_2,Zb_3,Zb_4)_N}{Z^N(q)_N}, & n~\text{odd},\\[4mm]
\displaystyle\frac{(Z,Zb_1b_2,Zb_1b_3,Zb_1b_4)_N}{(Zb_1)^N(q)_N}, & n~\text{even}.
\end{cases}\label{gri}
\end{gather}
\end{Theorem}

We will prove Theorem \ref{grt} by induction on $N$. In the case $N=1$, we have $x_i=\delta_{ik}$ for some~$k$. Using $k$ as summation index,
Theorem~\ref{grt} reduces to the following theta function identity.

\begin{Lemma}
For parameters subject to $b_1b_2b_3b_4z_1^2\dotsm z_n^2=1$,
\begin{gather}\label{ib}\sum_{k=1}^n\frac{\prod\limits_{j=1}^4\theta(z_kb_j)}{z_k}\prod\limits_{j=1,\,j\neq k}^n\frac{\theta(z_kz_j)}{\theta(z_k/z_j)}
=\begin{cases}
\theta(Zb_1,Zb_2,Zb_3,Zb_4)/Z, & n~\text{odd},\\
\theta(Z,Zb_1b_2,Zb_1b_3,Zb_1b_4)/Zb_1, & n~\text{even}.
 \end{cases} \end{gather}
\end{Lemma}

\begin{proof}
We apply induction on $n$, starting from the trivial case $n=1$. Let $b_1=vw$ and $b_2=v/w$, with $v$ and $w$ free parameters. As a function of $w$, each term in the sum, as well as the right-hand side, has the form $f(w)=C\theta(aw,a/w)$ with $C$ and $a$ independent of $w$. It is a~classical fact that any such function is determined by its values at two generic points. Indeed, Weierstrass' identity (which is equivalent to the case $n=2$ of~\eqref{tpf}) states that
\begin{gather*}f(w)=f(b)\frac{\theta(cw,c/w)}{\theta(cb,c/b)}+f(c)\frac{\theta(bw,b/w)}{\theta(bc,b/c)}, \end{gather*}
provided that $bc, b/c\notin p^{\mathbb Z}$. Thus, it suf\/f\/ices to verify~\eqref{ib} for two independent values of~$b_1$. Assuming $n\geq 2$, we choose $b_1=1/z_{n-1}$ and $b_1=1/z_n$. By symmetry, it is enough to consider the second case. Then, the term corresponding to $k=n$ cancels and we are reduced to an identity equivalent to \eqref{ib}, with $n$ replaced by $n-1$ and $b_1$ by $z_n$.
\end{proof}

We mention that it is not hard to deduce \eqref{ib} from classical theta function identities. Indeed, let $t=b_{n+1}$ in \eqref{tpf} (so that the right-hand side vanishes) and then make the substitutions \smash{$n\mapsto n+4$}, $(z_1,\dots,z_n)\mapsto (z_1,\dots,z_n,1,-1,\sqrt p,-1/\sqrt p)$, $(b_1,\dots,b_n)\mapsto(z_1^{-1},\dots,z_n^{-1},b_1^{-1},b_2^{-1}$, $b_3^{-1},b_4^{-1})$. Using the elementary identities $\theta(z^2)=\theta(z,-z,\sqrt pz,-\sqrt pz)$ and $2=\theta(-1,\sqrt p,-\sqrt p)$, one may deduce that the left-hand side of~\eqref{ib} is equal to
\begin{gather*}\frac 1 2\left({(-1)^{n+1}Z\prod\limits_{j=1}^4\theta(b_j)}+{Z\prod\limits_{j=1}^4\theta(-b_j)}+\frac 1{\sqrt p} {\prod\limits_{j=1}^4\theta\big(\sqrt pb_j\big)}
-\frac 1{\sqrt p} {\prod\limits_{j=1}^4\theta\big({-}\sqrt pb_j\big)}\right).
\end{gather*}
The fact that this equals the right-hand side of~\eqref{ib} follows from Jacobi's fundamental formulae \cite[Section~21.22]{ww}.

The inductive step in the proof of Theorem \ref{grt} is almost identical to that of \cite[Theorem~5.1]{r}. Denoting the right-hand side of~\eqref{gri} by
$R_N(Z;b_1,b_2,b_3,b_4)$ (where we for a moment consider~$Z$ as a free variable) we observe that, regardless of the parity of~$n$,
\begin{gather*}R_{N+1}(Z;b_1,b_2,b_3,b_4)=\frac{q^N\theta(q)}{\theta(q^{N+1})} R_1\big(q^NZ;q^{-N}b_1,b_2,b_3,b_4\big)R_N(Z;qb_1,b_2,b_3,b_4). \end{gather*}
Assuming \eqref{gri} for f\/ixed $N$, it follows that
\begin{gather*}R_{N+1}(Z;b_1,b_2,b_3,b_4)=\frac{q^N\theta(q)}{\theta(q^{N+1})} R_1\big(q^NZ;q^{-N}b_1,b_2,b_3,b_4\big)\\
\qquad{} \times \sum_{\substack{x_1,\dots,x_n\geq 0,\\x_1+\dots+x_n=N}}
\frac{\Delta(zq^x)}{\Delta(z)}\prod\limits_{1\leq i<j\leq n}q^{x_ix_j}(z_iz_j)_{x_i+x_j}\prod\limits_{i=1}^n
\frac{(qz_ib_1)_{x_i}\prod\limits_{j=2}^4(z_ib_j)_{x_i}}{z_i^{x_i}\prod\limits_{j=1}^n(qz_i/z_j)_{x_i}},
\end{gather*}
where $Z=z_1\dotsm z_n$ and $q^NBZ^2=1$. We pull the factor $R_1$ inside the sum and expand it using~\eqref{gri}, with~$z_i$ replaced by $q^{x_i}z_i$. This gives
\begin{gather*}R_{N+1}(Z;b_1,b_2,b_3,b_4)\\
\qquad{} =\frac{q^N\theta(q)}{\theta(q^{N+1})}\sum_{\substack{x_1,\dots,x_n\geq 0,\\x_1+\dots+x_n=N}}\sum_{\substack{y_1,\dots,y_n\geq 0,\\y_1+\dots+y_n=1}}
\frac{\Delta(zq^{x+y})}{\Delta(z)}\prod\limits_{1\leq i<j\leq n}q^{x_ix_j}(z_iz_j)_{x_i+x_j+y_i+y_j}\\
\qquad\quad{}\times\prod\limits_{i=1}^n
\frac{(qz_ib_1)_{x_i}\big(q^{x_i-N}z_ib_1\big)_{y_i}\prod\limits_{j=2}^4(z_ib_j)_{x_i+y_i}}{z_i^{x_i}(z_iq^{x_i})^{y_i}
\prod\limits_{j=1}^n(qz_i/z_j)_{x_i}\big(q^{1+x_i-x_j}z_i/z_j\big)_{y_i}}\\
\qquad{} =\frac{q^N\theta(q)}{\theta(q^{N+1})}\sum_{\substack{x_1,\dots,x_n\geq 0,\\x_1+\dots+x_n=N+1}}\sum_{\substack{y_1,\dots,y_n\geq 0,\\y_1+\dots+y_n=1}} \frac{\Delta(zq^{x})}{\Delta(z)}\prod\limits_{1\leq i<j\leq n}q^{x_ix_j-x_iy_j-x_jy_i}(z_iz_j)_{x_i+x_j}\\
\qquad\quad{}\times\prod\limits_{i=1}^n
\frac{(qz_ib_1)_{x_i-y_i}\big(q^{x_i-y_i-N}z_ib_1\big)_{y_i}\prod\limits_{j=2}^4(z_ib_j)_{x_i}}{z_i^{x_i}q^{(x_i-y_i)y_i}
\prod\limits_{j=1}^n(qz_i/z_j)_{x_i-y_i}\big(q^{1+x_i-x_j-y_i+y_j}z_i/z_j\big)_{y_i}},
\end{gather*}
where we replaced each $x_i$ by $x_i-y_i$ and used that $y_iy_j=0$ for $i\neq j$.

By elementary manipulations, using
\begin{gather*}\prod\limits_{i<j}q^{x_iy_j+x_jy_i}\prod\limits_i q^{(x_i-y_i)y_i}=q^{(x_1+\dots+x_n)(y_1+\dots+y_n)-(y_1^2+\dots+y_n^2)}=q^{(N+1)\cdot 1-1}=q^N, \end{gather*}
the expression above can be rewritten
\begin{gather*}R_{N+1}(Z;b_1,b_2,b_3,b_4)\\
=\frac{1}{\theta(q^{N+1})}\sum_{\substack{x_1,\dots,x_n\geq 0,\\x_1+\dots+x_n=N+1}}
\frac{\Delta(zq^{x})}{\Delta(z)}\prod\limits_{1\leq i<j\leq n}q^{x_ix_j}(z_iz_j)_{x_i+x_j}
\prod\limits_{i=1}^n\frac{(qz_ib_1)_{x_i}\prod\limits_{j=2}^4(z_ib_j)_{x_i}}{z_i^{x_i}\prod\limits_{j=1}^n(qz_i/z_j)_{x_i}}
\\
\qquad{} \times\sum_{\substack{y_1,\dots,y_n\geq 0,\\y_1+\dots+y_n=1}}\prod\limits_{i=1}^n
\frac{\big(q^{x_i-y_i-N}z_ib_1\big)_{y_i}\prod\limits_{j=1}^n\big(q^{1+x_i-y_i}z_i/z_j\big)_{y_i}}{\big(q^{1+x_i-y_i}z_ib_1\big)_{y_i}
\prod\limits_{j=1,\,j\neq i}^n\big(q^{x_i-x_j}z_i/z_j\big)_{y_i}}.
\end{gather*}
Writing $y_i=\delta_{ik}$, the inner sum takes the form
\begin{gather*}\sum_{k=1}^n\frac{\theta\big(q^{x_k-N-1}z_kb_1\big)\prod\limits_{j=1}^n\theta(q^{x_k}z_k/z_j)}{\theta(q^{x_k}z_kb_1)\prod\limits_{j=1,\,j\neq k}^n\theta(q^{x_k-x_j}z_k/z_j)}. \end{gather*}
By \eqref{tpf}, this can be evaluated as{\samepage
\begin{gather*}
\theta\big(q^{N+1}\big)\prod\limits_{i=1}^n\frac{\theta(z_ib_1)}{\theta(q^{x_i}z_ib_1)}=\theta\big(q^{N+1}\big)
\prod\limits_{i=1}^n\frac{(z_ib_1)_{x_i}}{(qz_ib_1)_{x_i}}\end{gather*}
and we arrive at \eqref{gri} with $N$ replaced by $N+1$. This completes the proof of Theorem~\ref{grt}.}
\begin{thebibliography}{99}
\bibitem[7]{gr}
Gustafson R.A., Rakha M.A., {$q$}-beta integrals and multivariate basic
 hypergeometric series associated to root systems of type~{$A_m$},
 \href{https://doi.org/10.1007/PL00001285}{\textit{Ann. Comb.}} \textbf{4} (2000), 347--373.

\bibitem[12]{r}
Rosengren H., Elliptic hypergeometric series on root systems, \href{https://doi.org/10.1016/S0001-8708(03)00071-9}{\textit{Adv.
 Math.}} \textbf{181} (2004), 417--447, \href{http://arxiv.org/abs/math.CA/0207046}{math.CA/0207046}.


\bibitem[16]{sp}
Spiridonov V.P., Theta hypergeometric integrals, \href{https://doi.org/10.1090/S1061-0022-04-00839-8}{\textit{St. Petersburg
 Math.~J.}} \textbf{15} (2003), 929--967, \href{http://arxiv.org/abs/math.CA/0303205}{math.CA/0303205}.

\bibitem[21]{ww}
Whittaker E.T., Watson G.N., A course of modern analysis, 4th ed., \href{https://doi.org/10.1017/CBO9780511608759}{\textit{Cambridge
 Mathematical Library}}, Cambridge University Press, Cambridge, 1996.

\end{thebibliography}
\end{document}
